%! Polynomial Equations, Inequalities, and Modeling — Unit 2, Lesson 2.7 — Lesson Plan
\documentclass[10pt]{article}
\usepackage{precalculus-article}
\usepackage{precalculus-boxes}
\newcommand{\TallMath}[1]{$\displaystyle #1\rule[-1.4em]{0pt}{3.2em}$}
\newcommand{\CourseName}{Precalculus}
\newcommand{\MeetingLength}{55 minutes}
\newcommand{\UnitNumberName}{Unit 2: Polynomial Functions \quad}
\newcommand{\LessonNumberName}{Lesson 2.7: Polynomial Equations, Inequalities, and Modeling}

\begin{document}
\begin{center}
    {\Large\bfseries \CourseName \\
    {\small\bfseries \UnitNumberName \LessonNumberName}}
\end{center}

\begin{tcolorbox}[colback=lilac, colframe=plum, boxrule=0.9pt, arc=2mm,
    left=2mm, right=2mm, top=1.5mm, bottom=1.5mm]
    \textbf{Primary Objective:} Students will solve polynomial equations by
    collecting every term on one side and applying the zero-product property;
    solve polynomial inequalities by sign analysis, using multiplicity to
    predict where the sign changes and reporting the answer in interval
    notation; construct a polynomial model from a physical context together
    with the domain restriction that context forces; and use the model --- and
    an inequality on it --- to answer a question in a sentence with units.
    \hfill\textit{\small (CED Topic 1.14 --- Skills 1.C, 3.B; plus
    calculus-prerequisite equation and inequality work)}
\end{tcolorbox}

\renewcommand{\arraystretch}{1.6}
\begin{skillbox}[Priority Ideas \& Skills]{goldbox}
\begin{tabularx}{\linewidth}{@{}X X@{}}
\textbf{AP Skills} &
\textbf{Key Understandings} \\
\midrule
\textbf{1.C} --- Construct new functions using given functions, useful in
modeling contexts. &
\begin{itemize}[leftmargin=1.2em,itemsep=2pt,topsep=0pt]
  \item A model can be built from the restrictions a mathematical or physical
        scenario imposes --- the algebra alone never supplies the domain.
        (1.14.A.1)
  \item A model can also be built by transforming a parent function to match a
        data set or a described scenario. (1.14.A.2)
  \item Technology fits a linear, quadratic, cubic, or quartic model to data by
        \textit{regression}; the analyst still chooses the degree. (1.14.A.3)
\end{itemize} \\
\textbf{3.B} --- Apply numerical results in a given mathematical or applied
context. &
\begin{itemize}[leftmargin=1.2em,itemsep=2pt,topsep=0pt]
  \item A model is used to draw conclusions about the scenario it came from:
        answer the question, then say what the number \textit{means}, in a
        sentence, with units. (1.14.C.1)
  \item A value rejected by the algebra can still be rejected a second time by
        the context --- a root outside the domain is not an answer.
\end{itemize} \\
\textbf{Calculus prerequisite} --- Solve polynomial equations and inequalities.
&
\begin{itemize}[leftmargin=1.2em,itemsep=2pt,topsep=0pt]
  \item A solution of $p(x) = 0$, a real zero of $p$, and an $x$-intercept of
        the graph are one object under three names (Lesson 2.3).
  \item Between consecutive real zeros the sign of $p$ cannot change, so one
        test value decides an entire interval.
  \item Odd multiplicity $\Rightarrow$ the sign changes at that zero; even
        multiplicity $\Rightarrow$ it does not (Lesson 2.3).
\end{itemize} \\
\end{tabularx}
\end{skillbox}

\begin{skillbox}[{Vocabulary, Concepts, \& Theorems}]{greenbox}
\renewcommand{\arraystretch}{1.4}
\begin{tabularx}{\linewidth}{@{} >{\leavevmode\bfseries\color{plum}}p{0.27\linewidth} X @{}}
Polynomial equation & An equation $p(x) = q(x)$ with polynomials on both sides; it is solved by rewriting it as $p(x) - q(x) = 0$. \\
Zero-product property & If a product of factors equals $0$, at least one factor equals $0$. It works only against $0$ --- never against any other number. \\
Polynomial inequality & A statement $p(x) > 0$, $\ge 0$, $< 0$, or $\le 0$; its solution is a union of intervals, not a list of numbers. \\
Sign chart & A number line cut at the real zeros of $p$, with the sign of $p$ recorded on each piece. \\
Test value & Any convenient input inside an interval; its sign is the sign of $p$ on the whole interval. \\
Domain restriction & The set of inputs a context allows. It comes from the scenario --- lengths are positive, a cut cannot exceed the sheet --- not from the formula. \\
Regression & A technology routine that fits a model of a chosen degree (linear, quadratic, cubic, quartic) to a data set. \\
\end{tabularx}
\end{skillbox}

\begin{skillbox}[Activate Prior Knowledge \& Spiral Review]{lilac}
    Please see attached warm-up (spiral review).
    Skills reviewed: factoring a polynomial completely by pulling out the
    greatest common factor and then factoring the remaining quadratic
    (Lesson 2.6), and reading real zeros with their multiplicities off a
    factored form and calling each one crossing or tangent (Lesson 2.3). Item 3
    evaluates $V(x) = x(20-2x)(30-2x)$ at $x = 2$ --- the students do not yet
    know what $V$ is, and that is deliberate: the same expression is built from
    scratch in Part 3.
\end{skillbox}

\begin{skillbox}[Hook]{lilac}
    \textbf{The card, the scissors, and the question.} Hold up a rectangular
    sheet of card, $20$ cm by $30$ cm. Cut a square of side $x$ from each of
    the four corners and fold the flaps up: an open box.
    \begin{enumerate}[itemsep=2pt]
        \item Ask: how big can $x$ be? Take answers until someone says
              ``$10$ --- past that the flaps meet and there is nothing left.''
              Write $0 < x < 10$ on the board and label it \textit{the context,
              not the algebra}.
        \item Project the pre-drawn graph of $V$ from the notes with the line
              $V = 1000$ drawn across it. Ask: which cuts give a box holding
              \textbf{at least} $1000$ cm\textsuperscript{3}?
    \end{enumerate}
    \textit{Discussion prompt:} that is not an equation --- it is an
    \textit{inequality}, and its answer is a stretch of the number line rather
    than a list. We need both halves of today's lesson to answer it, and we
    will, in Part 4.
\end{skillbox}

\begin{skillbox}[Lesson]{lilac}
\begin{multicols}{2}

\textbf{Part 1: Solving Polynomial Equations} (10 min)

\begin{itemize}
  \item Say the through-line of the unit out loud: a solution of $p(x) = 0$,
        a real zero of $p$, and an $x$-intercept of the graph are the same
        object under three names.
  \item The routine: (1) collect everything on one side so the other side is
        $0$; (2) factor completely; (3) set each factor to $0$.
  \item Work $x^3 + 6x = 5x^2$. Move everything left:
        $x^3 - 5x^2 + 6x = 0$, then $x(x-2)(x-3) = 0$, so $x = 0, 2, 3$.
  \item \textbf{Name the classic error before anyone makes it:} dividing both
        sides by $x$ at the start. It is legal only if $x \ne 0$ --- and
        $x = 0$ is one of the answers, so dividing deletes it.
\end{itemize}

\textbf{Part 2: Polynomial Inequalities by Sign Analysis} (15 min)

\begin{itemize}
  \item Keep the same polynomial: solve $x^3 - 5x^2 + 6x > 0$. The zeros
        $0, 2, 3$ cut the line into four intervals.
  \item The idea that makes it work: a polynomial can only change sign by
        passing through $0$, so \textit{inside} an interval with no zeros the
        sign cannot change. One test value settles the whole interval.
  \item Fill the sign chart together: $-1 \to -12$, $1 \to 2$,
        $2.5 \to -0.625$, $4 \to 8$. Answer $(0,2) \cup (3,\infty)$.
  \item Then the multiplicity connection back to 2.3: for
        $(x+2)(x-1)^2$, the sign changes at $-2$ (odd) and does \textit{not}
        change at $1$ (even). Solve $> 0$: the answer
        $(-2,1) \cup (1,\infty)$ has a hole punched in it at $1$ because the
        inequality is strict.
\end{itemize}

\columnbreak

\textbf{Part 3: Building the Model} (13 min)

\begin{itemize}
  \item Return to the card. Derive the three dimensions in front of them:
        height $x$, width $20 - 2x$, length $30 - 2x$ --- \textit{two} corners
        come off each side, which is where the $2x$ comes from.
  \item $V(x) = x(20-2x)(30-2x)$, and expanded
        $V(x) = 4x^3 - 100x^2 + 600x$. Ask which form is more useful for each
        job.
  \item \textbf{The AP move (1.14.A.1):} the domain. Ask for three separate
        reasons: $x > 0$, $20 - 2x > 0$, $30 - 2x > 0$. The binding one is
        $x < 10$. Write $0 < x < 10$ and say plainly that no step of the
        algebra produced it.
\end{itemize}

\textbf{Part 4: Answering Questions With the Model} (12 min)

\begin{itemize}
  \item Evaluate: $V(3) = 3(14)(24) = 1008$. Require the sentence: ``A
        $3$-cm cut gives a box holding $1008$ cubic centimeters.''
  \item Close the loop --- the hook question is an inequality on a model:
        $V(x) \ge 1000$. Reduce to $x^3 - 25x^2 + 150x - 250 \ge 0$, note
        $x = 5$ works, factor to $(x-5)(x^2 - 20x + 50)$, get
        $x = 10 \pm 5\sqrt{2}$.
  \item $10 + 5\sqrt{2} \approx 17.07$ is a perfectly good root and a
        perfectly bad answer --- it is outside $0 < x < 10$. Reject it by the
        context.
  \item Sign chart on the domain gives $10 - 5\sqrt{2} \le x \le 5$, about
        $2.93$ cm to $5$ cm. Finish with the sentence and the units.
  \item Mention regression (1.14.A.3) in one minute: given \textit{data}
        rather than a scenario, technology fits a cubic or quartic for you.
        You still choose the degree, and you still say what it means.
\end{itemize}
\end{multicols}
\end{skillbox}

\begin{skillbox}[Explicit Instruction --- The Sign-Analysis Routine]{lilac}
Model these five steps every time, and require them on the activity and
homework:
\begin{enumerate}[itemsep=2pt]
  \item Rewrite the inequality so one side is $0$.
  \item Factor completely and find every real zero.
  \item Mark the zeros on a number line; they cut it into intervals.
  \item Test one convenient value in each interval and record the sign.
  \item Read off the intervals that match the inequality. Include the zeros
        when the symbol is $\ge$ or $\le$; exclude them when it is $>$ or $<$.
        Report in interval notation, joining pieces with $\cup$.
\end{enumerate}
Worked live: $x^3 - 5x^2 + 6x > 0 \Rightarrow x(x-2)(x-3) > 0 \Rightarrow$
zeros $0, 2, 3 \Rightarrow$ signs $-, +, -, + \Rightarrow (0,2) \cup
(3,\infty)$.
\end{skillbox}

\begin{skillbox}[Active Monitoring]{redbox}
While students work on guided notes and practice, circulate and check:
\begin{itemize}
  \item \textbf{Common error:} dividing both sides by $x$ (or by $x - 2$) to
        ``simplify.'' Cold-call: ``What value of $x$ did you just assume was
        impossible? Check whether it is actually a solution.''
  \item \textbf{Common error:} setting factors equal to the number on the
        right, e.g.\ from $(x-1)(x+4) = 6$ writing $x - 1 = 6$. Push: ``The
        zero-product property is about \textit{zero}. Move the 6 first.''
  \item \textbf{Common error:} answering an inequality with a list of numbers
        instead of intervals. ``You solved the equation. The question asked
        where it is positive.''
  \item \textbf{Common error:} testing a value \textit{at} a zero instead of
        between zeros. Every test value must be strictly inside its interval.
  \item \textbf{Common error:} in the modeling half, giving the domain as
        ``all real numbers'' because the cubic is defined everywhere. Ask:
        ``Cut a $-3$ cm square. Show me.''
  \item \textbf{Language:} a bare number is not an answer to a modeling
        question. Require the noun and the unit every time.
\end{itemize}
\end{skillbox}

\begin{skillbox}[Group Work \& Differentiation]{redbox}
Students work on the \textbf{Group Activity: The Packaging Desk} in leveled
tiers. Every tier works the same physical context --- a $16$ cm $\times$ $24$
cm sheet --- so groups can be moved between tiers without new materials, and
no tier draws a graph.

\begin{itemize}
  \item \textbf{Tier R (Remediate):} Solve three factored or
        easily-factored equations, complete a two-interval sign chart for
        $(x+3)(x-2)$, and state the domain restriction for the sheet from the
        context.
  \item \textbf{Tier A (Approaching Proficiency):} Solve $x^3 = 2x^2 + 8x$
        with the everything-on-one-side move, complete a four-interval sign
        chart for $(x+2)(x-1)(x-3) \le 0$ and report in interval notation, then
        build $V(x) = x(16-2x)(24-2x)$ and evaluate $V(3)$ with a sentence.
  \item \textbf{Tier E (Extension):} Expand $V$ to $4x^3 - 80x^2 + 384x$,
        justify the domain $0 < x < 8$ from three separate restrictions, then
        solve $V(x) \ge 480$ --- factoring out $(x-2)$, solving
        $x^2 - 18x + 60 = 0$ for $9 \pm \sqrt{21}$, \textbf{rejecting
        $9 + \sqrt{21} \approx 13.58$ on domain grounds}, and reporting
        $[\,2,\ 9-\sqrt{21}\,]$ in a sentence with units (AP Skill 3.B).
\end{itemize}
\end{skillbox}

\begin{skillbox}[Individual Work \& Assessment]{redbox}
\textbf{Exit Ticket} (last 5--7 minutes, individual, collected):
\begin{enumerate}
  \item Solve $x^3 = 4x$ (worked space provided) --- the divide-by-$x$ trap.
  \item Complete a three-column sign chart for $p(x) = (x+3)(x-1)$ and give
        the solution of $p(x) < 0$ in interval notation.
  \item For a box cut from a $12$ cm square, $V(x) = x(12-2x)^2$: state the
        domain restriction the context forces, and compute $V(2)$ with units.
\end{enumerate}
\textbf{Assessment note:} Item 1 is the diagnostic --- an answer of
``$x = \pm 2$'' means the student divided by $x$ and lost a solution, and that
is a reteach, not a slip. On item 2 the target answer is $(-3,1)$; a student
who writes $\{-3, 1\}$ solved the equation instead of the inequality. Item 3(a)
is the 1.14.A.1 check: ``$0 < x < 6$'' earns it, ``all reals'' does not.
\end{skillbox}

\begin{skillbox}[Reinforcement \& Extension]{goldbox}
\textbf{Homework:} 6 problems --- two equations solved by factoring (one
requiring the everything-on-one-side move), a three-zero sign chart with the
answer in interval notation, a repeated-factor inequality where the sign fails
to change at the even-multiplicity zero, a full modeling item on an $18$ cm
$\times$ $24$ cm sheet (domain, expanded form, evaluation with units), a
regression item on choosing the degree from constant successive differences,
and one AP-style multiple-choice item on $x(x-2)^2 > 0$.

\textbf{Extension (optional):} Why is one test value per interval enough? A
polynomial is continuous, so to get from a positive output to a negative one
it must pass through $0$ --- and there are no zeros strictly inside the
interval. Students write the argument in two or three sentences; it is the
Intermediate Value Theorem in everything but name, and calculus will name it.

\textbf{Preview:} This closes Unit 2. The \textbf{Unit 2 Test} covers
polynomial behavior end to end --- end behavior, real and complex zeros,
multiplicity, equivalent forms, and today's equations, inequalities, and
models. Then \textbf{Unit 3, Rational Functions}, divides one polynomial by
another. Today's sign chart returns almost immediately there, with one new
rule: a rational expression can also change sign where its \textit{denominator}
is zero, at a point that is not in the domain at all.
\end{skillbox}

\begin{teachernote}[Warm-Up]
Five minutes, silent start, no calculators. Items 1 and 2 are straight recall
from 2.6 and 2.3 and should be quick --- do not reteach factoring here, just
confirm it, because every inequality today begins with a factoring step and a
class that is slow on item 1 will be slow all period. Item 3 is arithmetic on
purpose: students evaluate $V(2) = 832$ without being told what $V$ models.
Do not explain it during review. When the same expression is derived from the
card in Part 3, ask who recognizes it; the recognition is worth more than the
explanation would have been.
\end{teachernote}

\begin{teachernote}[Guided Notes]
Part 2 is the centerpiece and should get the most clock. Resist handing over
the sign chart as a procedure --- ask first why one test value is allowed to
speak for infinitely many, and let a student say ``because it would have to
cross zero to change sign.'' Write that sentence down; it is the homework
extension and it is calculus. Part 1 is short by design: the routine is three
steps and the only real content is the warning about dividing by $x$, which is
worth writing on the board in a different color. In Part 3 the derivation of
$20 - 2x$ deserves a physical demonstration --- fold an actual card, because
the students who write $20 - x$ are not making an algebra mistake, they are
picturing one corner instead of two. If time runs short, compress Part 1, never
Part 4: 1.14.C.1 lives entirely in Part 4 and it is where the two halves of the
lesson finally meet.
\end{teachernote}

\begin{teachernote}[Group Activity]
All three tiers share the $16 \times 24$ sheet, so a group can be moved up or
down mid-period with no new handout. Tier R is deliberately light on algebra:
the whole task is ``zeros cut the line, test one point per piece,'' repeated
until it is automatic. In Tier A, watch question 1 --- a group that starts by
dividing by $x$ will get two solutions instead of three, and that is the moment
to stop the room. Tier E question 3 is the payoff of the whole lesson: the
quadratic yields two roots and one of them, $9 + \sqrt{21} \approx 13.58$, is
mathematically flawless and physically absurd. Make a group say out loud
\textit{why} it is rejected; ``because the sheet is only 16 cm wide'' is the
answer, and ``because it is too big'' is not.
\end{teachernote}

\begin{teachernote}[Exit Ticket]
Collected, completion credit only. Sort by item 1 first: $\{0, 2, -2\}$ is the
target and $\{2, -2\}$ is the reteach pile, and the two piles will be roughly
the size of the split you saw while circulating. On item 2, accept $(-3,1)$
only --- brackets there are wrong because the inequality is strict, and it is
worth thirty seconds of whole-class correction tomorrow rather than a comment
on each paper. Item 3(b) should include the units; if more than a few papers
give a bare $128$, open the test review with the sentence requirement, since
the Unit 2 test scores it.
\end{teachernote}

\begin{teachernote}[Homework]
Problem 3 is the one to grade first --- it is the only item where a zero fails
to produce a sign change, and a student who reports $[-1,3] \cup [3,\infty)$
rather than $[-1,\infty)$ has the multiplicity idea but not the interval
notation, which is a very different conversation from one who reports
$[3,\infty)$ alone. Problem 6 is the AP-style item: distractor A, $(0,\infty)$,
is chosen by everyone who forgets that $x = 2$ makes the product zero and a
strict inequality excludes it, and it is worth two minutes of autopsy because
the same puncture appears on the unit test. Since this is the last homework of
the unit, collect the extension too --- the continuity argument is the one
piece of today that Unit 3 assumes and never restates.
\end{teachernote}
\end{document}
